\section{Vector Spaces}

\subsection{$\R^n$ and $\C^n$}

Before defining vector spaces, Axler establishes the fields we are working over. The symbol $\F$ will always denote either the real numbers $\R$ or the complex numbers $\C$.
\begin{definition}[Fields]
	A field is a set containing at least two distinct elements called 0 and 1, along with
	operations of addition and multiplication satisfying the properties of commutativity, associativity, identities, additive inverse, multiplicative inverse and distributivity.
\end{definition}
\subsubsection{Complex Numbers, $\C$}

\begin{definition}[Complex Numbers]
\leavevmode
\begin{itemize}
	\item A complex number is an ordered pair $(a,b)$ where $a,b \in \R$, but we will write this as $a + bi$.
	\item The set of all complex numbers is denoted by $\C$:
	\[ \C = \{a + bi \colon a, b \in \R\}. \]
	\item Addition and multiplication on $\C$ are defined by
	\[ (a + bi) + (c + di) = (a + c) + (b + d)i, \]
	\[ (a + bi)(c + di) = (ac - bd) + (ad + bc)i; \]
	here a, b, c, d $\in$ $\R$
\end{itemize}
\end{definition}
\begin{tcolorbox}[colback=yellow!20, colframe=yellow!80!black, title=My thoughts]
Just assume $i^2 = -1$ and everything works the same as in $\R$. Commutativity, associativity, identities, additive inverse, multiplicative inverse and distributivity are all properties of complex arithmetic.
\end{tcolorbox}
\subsubsection{Lists}
\begin{definition}[List, length]
\leavevmode
\begin{itemize}
	\item Suppose $n$ is a nonnegative integer. A list of length $n$ is an ordered collection of $n$ elements (which might be numbers, other lists, or more abstract
	objects).
	\item Two lists are equal if and only if they have the same length and the same
	elements in the same order.
\end{itemize}
\end{definition}
\begin{tcolorbox}[colback=yellow!20, colframe=yellow!80!black, title=My thoughts]
A list can have multiple same entries unlike sets. Fields and Vector Spaces (defined later) are sets, but a list is a tuple.
\end{tcolorbox}
\subsubsection{$\F^n$}
\begin{definition}[$\F^n$, coordinate]
	$\F^n$ is the set of all lists of length $n$ of elements of $\F$:
	\[
	\F^n = \{(x_1, \dots, x_n) : x_k \in \F \text{ for } k \in \{1, \dots, n\}\}
	\]
	For $(x_1, \dots, x_n) \in \F^n$ and $k \in \{1, \dots, n\}$, we say that $x_k$ is the $k^{th}$ coordinate of $(x_1, \dots, x_n)$.
\end{definition}
\begin{definition}[addition in $F^n$]
	Addition in $F^n$ is defined by adding corresponding coordinates:
	\[(x_1, … , x_n) + (y_1, … , y_n) = (x_1 + y_1, … , x_n + y_n)\]
\end{definition}
\begin{definition}[additive inverse in $\F^n$, $-x$]
	For $x \in \F^n$, the additive inverse of $x$, denoted by $-x$, is the vector $-x$ $\in$ $\F^n$
	such that
	\[x + (-x) = 0\]
	Thus if $x = (x_1, \dots, x_n)$, then $-x = (-x_1, \dots, -x_n)$.
\end{definition}
\begin{definition}[scalar multiplication in $\F^n$]
	The product of a number $\lambda$ and a vector in $\F^n$ is computed by multiplying
	each coordinate of the vector by $\lambda$:
	\[\lambda(x_1, \dots, x_n) = (\lambda x_1, \dots , \lambda x_n);\]
	here $\lambda$ $\in$ $\F$ and $(x_1, \dots , x_n) \in \F^n$
\end{definition}
\subsection{Definition of a Vector Space}

Now we abstract the properties of $\R^n$ and $\C^n$ into a general structure.

\begin{definition}[Vector Space]
	A vector space is a set $\V$ along with an addition on $\V$ and a scalar multiplication on $\V$ such that commutativity, associativity, identities, additive inverse, and distributivity hold.
\end{definition}
\begin{tcolorbox}[colback=yellow!20, colframe=yellow!80!black, title=My thoughts]
	Essentially, the difference between a field and a vector space is the fact that only scalar multiplication is defined for elements of a vector space. This also implies it is not necessary for a vector space to have multiplicative inverses. Also note, a field is a standalone set while a vector space relies on an underlying field to exist at all.
\end{tcolorbox}
\begin{definition}[Real Vector Space, Complex Vector Space]
\leavevmode
\begin{itemize}
	\item A vector space over $\R$ is called a real vector space.
	\item A vector space over $\C$ is called a complex vector space.
\end{itemize}
\end{definition}
\subsubsection{Some Property Proofs}
\begin{theorem}[Unique Additive Identity]
	A vector space has a unique additive identity.
\end{theorem}

\begin{proof}
	Suppose $0$ and $0'$ are both additive identities for $V$. Then
	\[ 0' = 0' + 0 = 0 + 0' = 0, \]
	where the first equality holds because 0 is an additive identity, the second equality
	comes from commutativity, and the third equality holds because $0'$ is an additive
	identity. Thus $0' = 0$, proving that $\V$ has only one additive identity.
\end{proof}
\begin{theorem}
	a number times the vector 0 is 0
\end{theorem}
\begin{proof}
	For $v \in \V$, we have
	$0v = (0 + 0)v = 0v + 0v$.
	Adding the additive inverse of $0v$ to both
	sides of the equation above gives $0 = 0v$,
	as desired.
\end{proof}
\begin{theorem}
	the number -1 times a vector gives the additive inverse of the vector
\end{theorem}
\begin{proof}
	For $v \in \V$, we have
	$v + (-1)v = 1v + (-1)v = (1 + (-1))v = 0v = 0$.
	This equation says that $(-1)v$, when added to $v$, gives 0. Thus $(-1)v$ is the
	additive inverse of $v$, as desired
\end{proof}
\subsection{Subspaces}
\begin{definition}[Subspace]
	A subset $U$ of $\V$ is called a subspace of $\V$ if $U$ is also a vector space under the same addition and scalar multiplication.
\end{definition}
\subsubsection{Sums of Subspaces}
\begin{definition}[Sum of Subspaces]
	Suppose $\V_1, \dots , \V_m$ are subspaces of $\V$. The sum of $\V_1, \dots , \V_m$, denoted by $\V_1 + \dots + \V_m$, is the set of all possible sums of elements of $\V_1,\dots , \V_m$. More
	precisely,
	\[\V_1 + \dots + \V_m = \{v_1 + \dots + v_m \colon v_1 \in \V_1, \dots , v_m \in \V_m\}\]
\end{definition}
\begin{theorem}
	Suppose $\V_1, \dots , \V_m$ are subspaces of $\V$. Then $\V_1 + \dots + \V_m$ is the smallest
	subspace of $\V$ containing $\V_1, \dots , \V_m$.
\end{theorem}
\begin{tcolorbox}[colback=yellow!20, colframe=yellow!80!black, title=My thoughts]
	I will add my own proof that sum of subspaces is a subspace (although Axler skips this) because I think it is necessary.
\end{tcolorbox}
\begin{proof}
	Let us first prove that sum of subspaces is a subspace
\leavevmode
\begin{itemize}
	\item $\V_1 + \dots + \V_m$ contains 0 (additive identity) since all the subspaces contain 0 and sum of subspaces contain every possible sum of their elements.
	\item Let $x, y \in (\V_1 + \dots + \V_m).$ \\
	By definition we can express them as $x = v_1 + \dots + v_m$ and $y = u_1 + \dots + u_m$ such that $u_1, v_1 \in \V_1, \dots, u_m, v_m \in \V_m$
	So we have \[x + y = (v_1 + \dots + v_m) + (u_1 + \dots + u_m)\]
	By commutativity and associativity of vector addition, \[x + y = (v_1 + u_1) + \dots + (v_m + u_m)\]
	Since $\V_1, \dots, \V_m$ are subspaces, $(v_1 + u_1) \in \V_1, \dots, (v_m + u_m) \in \V_m$
	
	Hence we conclude $(x + y) \in \V_1 + \dots + \V_m$
	
	Thus, it is closed under addition
	\item Let $x \in (\V_1 + \dots + \V_m).$ and $x = v_1 + \dots + v_m$ where $v_1 \in \V_1, \dots, v_m \in \V_m.$ Let $c$ be any scalar from the field $\F.$\\
	Multiplying $x$ with the scalar $c$ yields:\[cx = c(v_1 + \dots + v_m) = cv_1 + \dots + cv_m\]
	Since $\V_1 + \dots + \V_m$ are subspaces $cv_1 \in \V_1, \dots , cv_m \in \V_m.$\\
	We conclude $cx \in \V_1 + \dots + \V_m$
\end{itemize}
Thus, sum of subspaces is a subspace.

Now we can start with the proof for $\V_1, \dots , \V_m$ being the smallest subspace of $\V$ containing $\V_1, \dots , \V_m$.

The subspaces $\V_1, \dots , \V_m$ are all contained in $\V_1 + \dots + \V_m,$ considering sums $v_1 + \dots + v_m$ where all except one of the $v_k$’s are 0. 

Conversely, every subspace of $\V$ containing $\V_1, \dots , \V_m$ contains $\V_1 + \dots + \V_m$ (because subspaces
must contain all finite sums of their elements). Thus $\V_1 + \dots + \V_m$ is the smallest
subspace of $V$ containing $\V_1, \dots , \V_m.$
\end{proof}
\subsubsection{Direct Sums}
\begin{definition}[Direct Sum]
	Suppose $\V_1, \dots , \V_m$ are subspaces of $\V$.
\begin{itemize}
	\item The sum $\V_1 + \dots + \V_m$ is called a direct sum if each element of $\V_1 + \dots + \V_m$
	can be written in only one way as a sum $v_1 + \dots + v_m$, where each $v_k \in \V_k$.
	\item If $\V_1 + \dots + \V_m$ is a direct sum, then $\V_1 \oplus \dots \oplus \V_m$ denotes $\V_1 + \dots + \V_m$,
	with the $\oplus$ notation serving as an indication that this is a direct sum.
\end{itemize}	
\begin{theorem}
	Suppose $\V_1 + \dots + \V_m$ are subspaces of $\V$. Then $\V_1 + \dots + \V_m$ is a direct sum if
	and only if the only way to write 0 as a sum $v_1 + \dots + v_m$, where each $v_k \in \V_k$,
	is by taking each $v_k$ equal to 0.
\end{theorem}
\begin{proof}
	First suppose $\V_1 + \dots + \V_m$ is a direct sum. Then the definition of direct
	sum implies that the only way to write 0 as a sum $v_1 + \dots + v_m$, where each $v_k \in \V_k$,
	is by taking each $v_k$ equal to 0.
	
	\begin{tcolorbox}[colback=yellow!20, colframe=yellow!80!black, title=My thoughts]
		This is because taking every $v_k = 0$ is a trivial solution, but since we know it is a direct sum, the only solution would be taking every $v_k = 0$.
	\end{tcolorbox}
	
	Now suppose that the only way to write 0 as a sum $v_1 + \dots + v_m$, where each
	$v_k \in \V_k$, is by taking each $v_k$ equal to 0. To show that $\V_1 + \dots + \V_m$ is a direct
	sum, let $v \in \V_1 + \dots + \V_m$. We can write
	\[v = v_1 + \dots + v_m\]
	for some $v_1 \in \V_1, \dots, v_m \in \V_m$. To show that this representation is unique,
	suppose we also have
	\[v = u_1 + \dots + u_m\]
	where $u_1 \in \V_1, \dots , u_m \in \V_m$. Subtracting these two equations, we have
	\[0 = (v_1 - u_1) + \dots + (v_m - u_m)\]
	Because $v_1 - u_1 \in \V_1, \dots , v_m - u_m \in \V_m$, the equation above implies that each
	$v_k - u_k$ equals 0. Thus $v_1 = u_1, \dots , v_m = u_m$, as desired.
\end{proof}
\begin{theorem}
	Suppose $U$ and $W$ are subspaces of $V$. Then
	$$U + W \text{ is a direct sum } \iff U \cap W = \{0\}.$$
\end{theorem}
\begin{proof}
	First suppose that $U + W$ is a direct sum. If $v \in U \cap W$, then $0 = v +(-v)$,
	where $v \in U$ and $-v \in W$. By the unique representation of 0 as the sum of a
	vector in $U$ and a vector in $W$, we have $v = 0$. Thus $U \cap W = \{0\}$, completing
	the proof in one direction.
	To prove the other direction, now suppose $U \cap W = \{0\}$. To prove that $U + W$
	is a direct sum, suppose $u \in U, w \in W$, and
	$0 = u + w$.
	To complete the proof, we only need to show that $u = w = 0$. The
	equation above implies that $u = -w \in W$. Thus $u \in U \cap W$. Hence $u = 0$,
	which by the equation above implies that $w = 0$, completing the proof.
\end{proof}
\end{definition}
\subsection{Exercise 1A}
\begin{tcolorbox}[colback=yellow!20, colframe=yellow!80!black, title=My thoughts]
	I can't find a way to prove these identities by themselves, maybe I need to learn analysis for that. I will just be assuming $\R$ is already known to be a field, and use those properties to prove these.
\end{tcolorbox}
\begin{question}
	Show that $\alpha + \beta = \beta + \alpha$ for all $\alpha, \beta \in \C$.  
\end{question}
\begin{proof}
	Letting $\alpha = a+bi$, $\beta = c+di$, $\alpha + \beta = a + ib + c + id$
	\[= (a + c) + (b + d)i\]
	\[= (c + a) + (d + b)i\]
	\[= (c + id) + (a + ib)\]
	\[= \beta + \alpha\]
\end{proof}
\begin{question}
	Show that $(\alpha + \beta) + \lambda = \alpha + (\beta + \lambda)$ for all $\alpha, \beta, \lambda \in \C$.     
\end{question}
\begin{proof}
	Letting $\alpha = a+bi$, $\beta = c+di$, and $\lambda = e+fi$ (where $a, b, c, d, e, f \in \R$), the left side of the equation expands as:\\$(\alpha + \beta) + \lambda$\\$= ((a+c) + (b+d)i) + (e+fi)$\\$= ((a+c)+e) + ((b+d)+f)i$\\Because addition is associative in the real numbers, we can rewrite the real and imaginary components as $(a+(c+e))$ and $(b+(d+f))$.\\Expanding the right side gives:\\$\alpha + (\beta + \lambda)$\\$= (a+bi) + ((c+e) + (d+f)i)$\\$= (a+(c+e)) + (b+(d+f))i$
\end{proof}
\begin{question}
	Show that $(\alpha\beta)\lambda = \alpha(\beta\lambda)$ for all $\alpha, \beta, \lambda \in \C$.
\end{question}
\begin{proof}
	Letting $\alpha = a+bi$, $\beta = c+di$, and $\lambda = e+fi$ (where $a, b, c, d, e, f \in \R$), the left side of the equation expands as:\[((ac-bd) + (ad+bc)i)(e+fi)\] Multiplying this out yields a real part of $ace - bde - adf - bcf$ and an imaginary part of $acf - bdf + ade + bce$.\\Expanding the right side $\alpha(\beta\lambda)$ gives \[(a+bi)((ce-df) + (cf+de)i)\] Multiplying this out yields a real part of $ace - adf - bcf - bde$ and an imaginary part of $acf + ade + bce - bdf$.
\end{proof}
\begin{question}
	Show that $\lambda(\alpha + \beta) = \lambda\alpha + \lambda\beta$ for all $\lambda, \alpha, \beta \in \C$
\end{question}
\begin{proof}
	Expanding $\lambda(\alpha + \beta)$ with $\lambda = e+fi$, $\alpha = a+bi$, and $\beta = c+di$ gives a real part of $ea + ec - fb - fd$ and an imaginary part of $eb + ed + fa + fc$. Multiplying out $\lambda\alpha + \lambda\beta$ yields the exact same terms, relying on the distributive property of real numbers to match them up.
\end{proof}
\begin{question}
	Show that for every $\alpha \in \C$, there exists a unique $\beta \in \C$ such that $\alpha + \beta = 0$. 
\end{question}
\begin{proof}
	Setting $\alpha + \beta = 0$ with $\alpha = a+bi$ and $\beta = c+di$ results in $(a+c) + (b+d)i = 0$.\\ Equating the real and imaginary components to zero means $a+c = 0$ and $b+d = 0$, yielding exactly $c = -a$ and $d = -b$. Because every real number has a uniquely defined additive inverse in $\R$, the components $c$ and $d$ exist and are unique, making $\beta$ unique.
\end{proof}
\begin{question}
	Show that for every $\alpha \in \C$ with $\alpha \neq 0$, there exists a unique $\beta \in \C$ such that $\alpha\beta = 1$.
\end{question}
\begin{proof}
	$\alpha = a+bi$, $\beta = c+di$\\
	Setting $(a+bi)(c+di) = 1$ results in $(ac - bd) + (ad + bc)i = 1 + 0i$. Equating the real and imaginary parts gives a system of two linear equations: $ac - bd = 1$ and $ad + bc = 0$. Since $\alpha \neq 0$, we know that $a$ and $b$ are not both zero, meaning $a^2 + b^2 \neq 0$. Solving this system for $c$ and $d$ yields exactly $c = \frac{a}{a^2+b^2}$ and $d = \frac{-b}{a^2+b^2}$. This produces a single, well-defined real value for both $c$ and $d$.
\end{proof}
\begin{question}
	Show that $\frac{-1 + \sqrt{3}i}{2}$ is a cube root of 1 (meaning that its cube equals 1).
\end{question}
\begin{proof}
	Squaring $\frac{-1 + \sqrt{3}i}{2}$ yields  $\frac{-1 - \sqrt{3}i}{2}$. Multiplying that squared result by the original term $\frac{-1 + \sqrt{3}i}{2}$ gives $\frac{(-1)^2 - (\sqrt{3}i)^2}{4} = \frac{1 - (-3)}{4}$, which reduces perfectly to $1$.
\end{proof}
\begin{question}
	Find two distinct square roots of $i$.
\end{question}
\begin{proof}
	$\frac{1}{\sqrt{2}} + \frac{1}{\sqrt{2}}i$ and $-\frac{1}{\sqrt{2}} - \frac{1}{\sqrt{2}}i$
\end{proof}
\begin{question}
	Find $x \in \R^4$ such that $(4, -3, 1, 7) + 2x = (5, 9, -6, 8)$.
\end{question}
\begin{proof}
	$(1/2, 6, -7/2, 1/2)$
\end{proof}
\begin{question}
	Explain why there does not exist $\lambda \in \C$ such that $\lambda(2 - 3i, 5 + 4i, -6 + 7i) = (12 - 5i, 7 + 22i, -32 - 9i)$. 
\end{question}
\begin{proof}
	Setting $\lambda(2 - 3i) = 12 - 5i$ for the first component gives $\lambda = \frac{12 - 5i}{2 - 3i}$.\\ Multiplying the numerator and denominator by the conjugate $(2 + 3i)$ yields $\frac{24 + 36i - 10i + 15}{13}$, which simplifies to $\lambda = 3 + 2i$.\\Testing this $\lambda$ on the second component yields $(3 + 2i)(5 + 4i) = 15 + 12i + 10i - 8 = 7 + 22i$, which matches the required vector perfectly.\\Now, applying $\lambda = 3 + 2i$ to the third component gives $(3 + 2i)(-6 + 7i) = -18 + 21i - 12i - 14 = -32 + 9i$. But the required third component in the target vector is $-32 - 9i$.
\end{proof}
\begin{question}
	Show that $(x + y) + z = x + (y + z)$ for all $x, y, z \in \F^n$.   
\end{question}
\begin{proof}
	Let $x = (x_1, \dots, x_n)$, $y = (y_1, \dots, y_n)$, and $z = (z_1, \dots, z_n)$. Adding component-wise, the left side gives $((x_1+y_1)+z_1, \dots, (x_n+y_n)+z_n)$. Because addition is associative in the underlying field $\F$, each component $(x_i+y_i)+z_i$ is equal to $x_i+(y_i+z_i)$. This perfectly constructs the right side $x+(y+z)$.
\end{proof}
\begin{question}
	Show that $(ab)x = a(bx)$ for all $x \in \F^n$ and all $a, b \in \F$.   
\end{question}
\begin{proof}
	Expanding $x$ as $(x_1, \dots, x_n)$ means the left side $(ab)x$ becomes $((ab)x_1, \dots, (ab)x_n)$. The right side $a(bx)$ evaluates to $a(bx_1, \dots, bx_n)$, which then scales to $(a(bx_1), \dots, a(bx_n))$. Because multiplication within the underlying field $\F$ is associative, $(ab)x_i = a(bx_i)$ for every component
\end{proof}
\begin{question}
	Show that $1x = x$ for all $x \in \F^n$.  
\end{question}
\begin{proof}
	Expanding $x$ into its components $(x_1, \dots, x_n)$ gives $1x = (1x_1, \dots, 1x_n)$. Because $1$ is the multiplicative identity in the field $\F$, $1x_i = x_i$ for every individual component. This gives us the original vector $x$.
\end{proof}
\begin{question}
	Show that $\lambda(x + y) = \lambda x + \lambda y$ for all $\lambda \in \F$ and all $x, y \in \F^n$.   
\end{question}
\begin{proof}
	Expanding $\lambda(x + y)$ component-wise gives $(\lambda(x_1 + y_1), \dots, \lambda(x_n + y_n))$. Because the underlying field $\F$ has the distributive property, each component expands to $\lambda x_i + \lambda y_i$. This perfectly matches the component-wise expansion of the right side, $\lambda x + \lambda y = (\lambda x_1 + \lambda y_1, \dots, \lambda x_n + \lambda y_n)$.
\end{proof}
\begin{question}
	Show that $(a + b)x = ax + bx$ for all $a, b \in \F$ and all $x \in \F^n$.   
\end{question}
\begin{proof}
	Expanding $(a + b)x$ component-wise yields $((a + b)x_1, \dots, (a + b)x_n)$. Applying the distributive property of the underlying field $\F$ to each component gives $(ax_1 + bx_1, \dots, ax_n + bx_n)$. This perfectly matches the right side $ax + bx$ when you add the scaled vectors $(ax_1, \dots, ax_n)$ and $(bx_1, \dots, bx_n)$ together.
\end{proof}
\subsection{Exercise 1B}
\setcounter{question}{0}
\begin{question}
	Prove that $-(-v) = v$ for every $v \in \V$.   
\end{question}
\begin{proof}
	$$-(-v) = -(-v) + 0$$$$-(-v) = -(-v) + (-v + v)$$$$-(-v) = (-(-v) + -v) + v$$$$-(-v) = 0 + v$$$$-(-v) = v$$
\end{proof}
\begin{question}
	Suppose $a \in \F, v \in \V,$ and $av = 0$. Prove that $a = 0$ or $v = 0$.   
\end{question}
\begin{proof}
	Given equation:$$av = 0$$
	Since $a \in \F$ (a field) and $a \neq 0$, it has a multiplicative inverse $a^{-1} \in \F$.
	$$a^{-1}(av) = a^{-1}(0)$$
	$$a^{-1}(av) = 0$$
	By associativity of scalar multiplication:
	$$(a^{-1}a)v = 0$$
	$$1v = 0$$
	$$v = 0$$
	If $a \neq 0$, then $v$ must be $0$. Therefore, if $av = 0$, it must be true that either $a = 0$ or $v = 0$.
\end{proof}
\begin{question}
	Suppose $v, w \in \V$. Explain why there exists a unique $x \in \V$ such that $v + 3x = w$.   
\end{question}
\begin{proof}
	First we prove existence:
	$$x = \frac{1}{3}(w - v)$$
	Since a vector space is closed under addition and scalar multiplication $x \in \V$\\
	Now for uniqueness:\\
	Suppose there are two vectors, $x_1$ and $x_2$ in $V$, that both satisfy the equation.
	$$v + 3x_1 = v + 3x_2$$
	$$-v + (v + 3x_1) = -v + (v + 3x_2)$$
	Using associativity and additive inverse properties,
	$$3x_1 = 3x_2$$
	$$x_1 = x_2$$
	Thus x is unique.
\end{proof}
\begin{question}
	The empty set is not a vector space. The empty set fails to satisfy only one of the requirements listed in the definition of a vector space. Which one?   
\end{question}
\begin{proof}
	It doesn't have the zero element.
\end{proof}
\begin{question}
	Show that in the definition of a vector space, the additive inverse condition can be replaced with the condition that $0v = 0$ for all $v \in \V$.   
\end{question}
\begin{proof}
	Let's assume the additive inverse condition does not exist. Instead we have $0v = 0$ for all $v \in \V$. We can arrive at the additive inverse criterion from here.
	$$0v = 0$$
	$$(1+(-1))v = 0$$
	Using distributivity,
	$$1v + (-1)v = 0$$
	Using multiplicative identity,
	$$v + (-1)v = 0$$
	The element $(-1)v$ satisfies the exact requirements of an additive inverse.
\end{proof}
\begin{question}
	Let $\infty$ and $-\infty$ denote two distinct objects, neither of which is in $\R$. Define an addition and scalar multiplication on $\R \cup \{\infty, -\infty\}$ as you could guess from the notation. Specifically, the sum and product of two real numbers is as usual, and for $t \in \R$ define
	\[
	t\infty = \begin{cases} 
		-\infty & \text{if } t < 0, \\
		0 & \text{if } t = 0, \\
		\infty & \text{if } t > 0,
	\end{cases}
	\qquad
	t(-\infty) = \begin{cases} 
		\infty & \text{if } t < 0, \\
		0 & \text{if } t = 0, \\
		-\infty & \text{if } t > 0,
	\end{cases}
	\]
	and
	\begin{align*}
		t + \infty &= \infty + t = \infty + \infty = \infty, \\
		t + (-\infty) &= (-\infty) + t = (-\infty) + (-\infty) = -\infty, \\
		\infty + (-\infty) &= (-\infty) + \infty = 0.
	\end{align*}
	With these operations of addition and scalar multiplication, is $\R \cup \{\infty, -\infty\}$ a vector space over $\R$? Explain.
\end{question}
\begin{proof}
	It is not a vector space.\\
	If t + $\infty = \infty$, we have:
	$$1 + \infty = 2 + \infty$$
	$$1=2$$
	We can cancel out the $\infty$ on both sides due to existence of additive inverse.
\end{proof}
\begin{question}
	Suppose $S$ is a nonempty set. Let $\V^S$ denote the set of functions from $S$ to $\V$. Define a natural addition and scalar multiplication on $\V^S$, and show that $\V^S$ is a vector space with these definitions.   
\end{question}
\begin{proof}
	Because $\V$ is a vector space, $\V^S$ will obey all vector space axioms through point-wise evaluations for every $x \in S$.
	\begin{itemize}
	\item Commutativity \& Associativity: $(f+g)(x) = f(x) + g(x) = g(x) + f(x) = (g+f)(x)$. Associativity follows the exact same logic.
	\item Additive Identity: The zero function, defined by $0(x) = 0$ for all $x \in S$.
	\item Additive Inverse: The function defined point-wise by $(-f)(x) = -(f(x))$.
	\item Distributive Properties: $a(f+g)(x) = a(f(x) + g(x)) = af(x) + ag(x) = (af + ag)(x)$.
\end{itemize}
\end{proof}
\begin{question}
	\textbf{8}\quad Suppose $\V$ is a real vector space.
	\begin{itemize}
		\item The \textit{complexification} of $\V$, denoted by $\V_{\C}$, equals $\V \times \V$. An element of $\V_{\C}$ is an ordered pair $(u, v)$, where $u, v \in \V$, but we write this as $u + iv$.
		\item Addition on $\V_{\C}$ is defined by
		\[
		(u_1 + iv_1) + (u_2 + iv_2) = (u_1 + u_2) + i(v_1 + v_2)
		\]
		for all $u_1, v_1, u_2, v_2 \in \V$.
		\item Complex scalar multiplication on $V_{\C}$ is defined by
		\[
		(a + bi)(u + iv) = (au - bv) + i(av + bu)
		\]
		for all $a, b \in \R$ and all $u, v \in \V$.
	\end{itemize}
	
	Prove that with the definitions of addition and scalar multiplication as above, $\V_{\C}$ is a complex vector space.
\end{question}
\begin{proof}
	Because $\V$ is a real vector space, it is inherently closed under its own addition and scalar multiplication. Therefore, $(u_1 + u_2) \in \V$, $(v_1 + v_2) \in \V$, $(au - bv) \in V$, and $(av + bu) \in \V$, guaranteeing the operations always produce valid elements of $\V_\C$. \\
	The following properties are also satisfied:\\
	\begin{itemize}
	\item Additive Identity: $0 + i0$ (since $0 \in \V$).
	\item Additive Inverse: $-(u) + i(-v)$ (since $\V$ contains inverses for $u$ and $v$).
	\item Multiplicative Identity: $1 + i0$ acts as the scalar $1$, yielding $1u - 0v + i(1v + 0u) = u + iv$.
	\item Associativity, Commutativity, and Distributivity: These expand algebraically and hold true because the vectors in $\V$ and scalars in $\R$ already obey these rules.
	\end{itemize}
\end{proof}